\chapter*{Introduction générale}
\addcontentsline{toc}{chapter}{Introduction générale} % to include
Dans un contexte où la santé et le bien-être occupent une place centrale dans nos vies quotidiennes, la gestion des médicaments représente un défi majeur pour de nombreuses familles et équipes soignantes. Entre les traitements multiples, les dates d'expiration à surveiller, les posologies à respecter et la coordination entre plusieurs membres d'une même famille, il devient difficile d'assurer un suivi rigoureux et sécurisé des médicaments.

Les méthodes traditionnelles de gestion, basées sur des carnets papier, des boîtes à pharmacie désorganisées ou des rappels manuscrits, présentent de nombreuses limites. Elles exposent les utilisateurs à des risques tels que l'oubli de prise de médicaments, la consommation de produits périmés, le manque de visibilité sur les stocks disponibles, ou encore l'absence de coordination entre les membres d'une famille ou d'une équipe de soins. Ces lacunes peuvent avoir des conséquences importantes sur la santé des patients et la qualité de leur prise en charge.

Face à ces problématiques, la transformation numérique offre aujourd'hui des opportunités inédites pour repenser la manière dont nous gérons nos médicaments au quotidien. L'utilisation d'applications mobiles permet désormais de centraliser l'information, d'automatiser les rappels et de faciliter la collaboration entre proches.

C'est dans ce contexte que s'inscrit notre projet de fin d'études, qui consiste à développer AFYA, une solution mobile de gestion intelligente des médicaments. Cette application vise à accompagner les utilisateurs dans l'organisation de leur pharmacie familiale, le suivi de leurs traitements et la coordination des soins au sein d'un groupe (famille, équipe soignante). L'objectif principal est de simplifier la gestion quotidienne des médicaments, d'améliorer l'observance thérapeutique et de renforcer la sécurité sanitaire des utilisateurs.

Au-delà de la gestion individuelle, ce projet se veut également un outil de collaboration familiale, en permettant aux membres d'un même foyer ou d'une équipe de soins de partager les informations médicales, de gérer collectivement la pharmacie et d'assurer un meilleur suivi des traitements de leurs proches.  

Ce rapport retrace les différentes étapes de réalisation de ce projet, structurées autour de cinq chapitres principaux :

\begin{itemize}
  \item Dans un premier chapitre, nous présenterons le \textbf{contexte général et les objectifs du projet}.
  Ce chapitre commencera par une présentation de l'organisme d'accueil, en mettant en lumière ses activités et les problématiques identifiées dans le domaine de la gestion des médicaments.
  Ensuite, nous réaliserons une \textit{étude de l'existant}, permettant d'analyser les solutions actuellement disponibles sur le marché ainsi que leurs limites.
  À partir de cette analyse, nous formulerons une proposition de solution adaptée aux besoins identifiés.
  Enfin, nous détaillerons la méthodologie adoptée pour la réalisation du projet, structurée selon une approche itérative et incrémentale permettant une livraison progressive de valeur.

  \item Le deuxième chapitre sera consacré à l'\textbf{analyse et la spécification des besoins}.
  Nous commencerons par identifier les différents acteurs du système et préciser leurs rôles respectifs (utilisateur, gestionnaire de groupe, responsable de soins).
  Ensuite, nous formaliserons les besoins fonctionnels et non fonctionnels de l'application.
  Nous établirons un diagramme de cas d'utilisation global afin de représenter graphiquement les interactions entre les utilisateurs et l'application.
  Ce chapitre présentera également le \textit{backlog produit}, regroupant l'ensemble des fonctionnalités planifiées, ainsi que la conception générale du système et les outils retenus pour sa réalisation.

  \item Le troisième chapitre portera sur la \textbf{Release 1 : Mise en place des fonctionnalités de base}.
  Ce chapitre décrira en détail la première étape du développement, dédiée aux fonctionnalités essentielles qui constituent le socle de l'application.
  Nous y présenterons les modules relatifs à la création de compte et à l'accès sécurisé, à la gestion des profils utilisateurs, ainsi qu'à la création et la gestion des groupes familiaux ou médicaux.
  Cette release introduit également les premières fonctionnalités de gestion des médicaments, permettant d'ajouter des produits à sa pharmacie personnelle et de consulter les informations associées.
  L'organisation du travail sera retracée à travers les planifications des phases concernées.
  Enfin, nous illustrerons cette version par des diagrammes conceptuels, les cas d'utilisation détaillés et un aperçu des écrans développés.

  \item Le quatrième chapitre présentera la \textbf{Release 2 : Gestion avancée des traitements et rappels}.
  Cette phase enrichit l'application avec des fonctionnalités permettant de créer et suivre des traitements médicaux complets, d'enregistrer des prescriptions et de configurer des rappels de prise de médicaments.
  Nous détaillerons les mécanismes de suivi des prises, de gestion des renouvellements et d'alerte en cas de stock faible ou de date d'expiration proche.
  Cette version met l'accent sur l'amélioration de l'observance thérapeutique grâce à des notifications intelligentes et un suivi détaillé de l'historique des prises.
  Des exemples concrets de scénarios d'utilisation seront présentés, accompagnés de diagrammes et d'interfaces illustrant les fonctionnalités développées.

  \item Enfin, un cinquième et dernier chapitre sera dédié à la \textbf{Release 3 : Fonctionnalités collaboratives et de gestion familiale}.
  Dans cette étape, nous nous concentrerons sur l'enrichissement de l'application par des fonctionnalités facilitant la collaboration au sein des groupes.
  Nous y aborderons notamment la gestion des membres du groupe, l'attribution de rôles et de responsabilités, ainsi que le suivi des données médicales des personnes à charge.
  Ce chapitre exposera également les fonctionnalités permettant de partager l'historique des traitements, de gérer collectivement la pharmacie familiale et de coordonner les soins entre plusieurs utilisateurs.
  Cette version vient ainsi compléter le projet en renforçant sa dimension collaborative et son utilité pour les familles et les équipes de soins.
\end{itemize}



\markboth{Introduction générale}{} %To redefine the section page head
